\documentclass{article}
\usepackage{graphicx}


\title{CW 

Gnu Linux}

\date{12/1/2025}

\author{Ehsan Esmaeili}

\begin{document}
\maketitle
\newpage
\section*{GNU's advantage:}

\begin{itemize}
    \item GNU is a free software project aiming to 
    create a completely free operating system without 
    reliance on proprietary software.
    \item Free software means the freedom 
    to run, study, modify, and share the software, 
    opposing restrictive proprietary licenses.
    \item The GNU General Public License (GPL) was developed 
    to ensure software freedom for users and developers.
\end{itemize}



\section*{Kernal:}

\subsection*{Kernal generally}

The kernel is the core of an operating system 
that manages communication between software and hardware. 
Its main tasks include managing hardware 
resources, processes, memory, and devices.
\subsection*{Kernal module}

Kernel modules are extensions that enhance 
the functionality of the kernel without requiring 
a full rebuild. They are used to add device 
drivers or specialized features.

\section*{Systemd:}

\begin{itemize}
    \item Systemd is an init system responsible for 
    managing the system boot process and services. 
    It handles background processes, daemons, 
    and system stability.
    \item Some criticize Systemd for violating 
    the Unix philosophy of creating small, 
    modular tools for specific tasks, arguing 
    that it adds unnecessary complexity to the system.
    \item Systemd services are processes running in the 
    background to handle tasks like network management, 
    task scheduling, and system functionalities.

    
\end{itemize}

\section*{Graphicx:}

\subsection*{Display server protocol}

A display server protocol is an interface between 
the operating system and the graphics hardware to 
render visuals on the screen. Examples include X11 and Wayland.
\subsection*{Desktop environment or windows man-ager}

A window manager controls the appearance and behavior 
of windows, while a desktop environment is a 
complete user interface that includes 
a window manager, menus, tools, and settings for 
a full experience.
\subsection*{Some examples}

Window managers: i3, Openbox, Fluxbox
Desktop environments: GNOME, KDE Plasma, XFCE




\section*{Ricing:}

\subsection*{Meaning of ricing}


Ricing refers to customizing 
the graphical appearance of a 
Linux system to make it visually appealing. 
It involves tweaking themes, icons, 
window manager configurations.
\subsection*{Favorite ricing}

\begin{itemize}
    \item Openbox
    \item XFCE 
\end{itemize}

\section*{Suckless software:}

\subsection*{Explane of Suckless software}

Suckless Software is a collection of 
lightweight and minimalist tools designed 
for simplicity, speed, and usability by advanced users.
\subsection*{Philosophy of software}

Their philosophy revolves around minimalism, 
adhering to the principle of "less is more" 
to make code readable, simple, and user-modifiable.

\section*{Shells:}

\subsection*{POSIX-compliant shell}

A POSIX-compliant shell, like Bash or Zsh, 
adheres to POSIX standards for ensuring 
compatibility and consistent behavior across v
arious operating systems.


\subsection*{Use full command:}


\subsubsection*{(A):}

grep -r "hello" testDir | wc -l


\subsubsection*{(B):}

ls -l  allthefiles.txt

\subsubsection*{(C):}

find testDir -type f -name "*.c" | wc -l


\end{document}
